% 31_body.tex — 산출물 ㉛ 코호트 리텐션 판독서 (빈 템플릿 / 완성 예시 공용 본문) — gen_product_templates.py 가 Product_specs.py 에서 생성
% 드라이버: 31_코호트판독서_빈템플릿.tex (\wbblanktrue) / 31_코호트판독서_완성예시.tex (\wbblankfalse — 워크북 §X.5 확정 후)
\documentclass[10pt,a4paper]{article}
\usepackage[a4paper,margin=16mm,top=14mm,bottom=20mm]{geometry}
\def\DocNum{31}\def\DocDay{7}\def\DocIdx{1}
\def\DocTitle{코호트 리텐션 판독서}
\def\DocSub{Cohort Retention Reading \& Allocation Decision}
\def\DocVariant{\B{빈 템플릿}{딥게이지 완성 예시}}
\usepackage{wbstyle}
\def\DocCirc{㉛}   % newunicodechar(㉛~㊵ 폴백) 뒤에 정의해야 활성 문자로 읽힌다
\begin{document}
\docheader
 {\Bg{[회사명 · 제품명]}{딥게이지 · GAUGE}}
 {\Bg{[v1.x / D+600 2027-10-23]}{v1.0 / D+600 2027-10-23}}
 {\Bg{[작성자 — CPO·CS]}{윤하경 CPO · 강지우 CS\rd{7mm}}}
 {\Bg{[검토자 — CEO]}{강민수\rd{7mm}}}

%% ------------------------------------------------------------------ 1
\secbar{1. 데이터 정의 — 분모를 먼저}
\guide{코호트 단위(월/분기) · **분모 정의**(신규 유료 계약 로고 / 활성 계정 / …) · 기준일 · 관측 가능한 셀(미도달 셀은 ▨). 이 표가 비면 아래 표는 읽을 수 없다.}
{\footnotesize
\begin{tabularx}{\linewidth}{|L{40mm}|Y|Y|Y|}\hline
\hd{항목} & \hd{자사(Ret-C)} & \hd{익명 A(Ret-A)} & \hd{익명 B(Ret-B)} \\ \hline
\ifwbblank
\rd{10mm}코호트 단위 &  &  &  \\ \hline
\rd{10mm}분모 정의 &  &  &  \\ \hline
\rd{10mm}분모 고정 시점 &  &  &  \\ \hline
\rd{10mm}기준일 &  &  &  \\ \hline
\rd{10mm}미도달 셀 &  &  &  \\ \hline
\else
코호트 단위 & 분기(계약일 기준) & 분기 & 분기 \\ \hline
분모 정의 & 해당 분기 신규 유료 계약 로고 & 확인 결과: 계약 후 첫 달 말 30일 내 로그인 로고 = 활성 계정 — 미로그인 로고 제외 & 자사와 동일(확인) \\ \hline
분모 고정 시점 & 계약일 — 이후 해지해도 분모 유지 & 첫 달 말(M0·M1 이탈은 보이지 않는다) & 계약일 \\ \hline
기준일 & 2027-10-23 & 2027-10-23(인쇄 기준) & 동일 \\ \hline
미도달 셀 & 27Q2 M5·M6 ▨ & 동일 & 동일 \\ \hline
\fi
\end{tabularx}}

%% ------------------------------------------------------------------ 2
\landscapeon
\secbar{2. 삼각행렬 — 자사 Ret-C (분자/분모 · \%)}
\guide{M0~M6. 분자·분모를 함께 적는다(정수여야 한다). 정본 §7.2.}
{\footnotesize
\begin{tabularx}{\linewidth}{|L{22mm}|L{14mm}|L{16mm}|L{16mm}|L{16mm}|L{16mm}|L{16mm}|L{16mm}|L{16mm}|L{22mm}|}\hline
\hd{코호트} & \hd{분모} & \hd{M0} & \hd{M1} & \hd{M2} & \hd{M3} & \hd{M4} & \hd{M5} & \hd{M6} & \hd{마지막 관측} \\ \hline
\ifwbblank
\rd{10mm}2026-Q4 &  &  &  &  &  &  &  &  &  \\ \hline
\rd{10mm}2027-Q1 &  &  &  &  &  &  &  &  &  \\ \hline
\rd{10mm}2027-Q2 &  &  &  &  &  &  &  &  &  \\ \hline
\else
2026-Q4 & 3 & 3 · 100 & 3 · 100 & 3 · 100 & 3 · 100 & 3 · 100 & 2 · 67 & 2 · 67 & M6 67\%(2/3) \\ \hline
2027-Q1 & 14 & 14 · 100 & 14 · 100 & 13 · 93 & 12 · 86 & 11 · 79 & 10 · 71 & 10 · 71 & M6 71\%(10/14) \\ \hline
2027-Q2 & 12 & 12 · 100 & 12 · 100 & 11 · 92 & 11 · 92 & 10 · 83 & ▨ & ▨ & M4 83\%(10/12) \\ \hline
\fi
\end{tabularx}}
\landscapeoff

%% ------------------------------------------------------------------ 3
\secbar{3. 평탄화 판정}
\guide{평탄선 높이 · 도달 시점 · 동일 연령(M4) 비교(코호트가 개선되는가) · '평탄'인가 '아직 떨어지는 중'인가.}
{\footnotesize
\begin{tabularx}{\linewidth}{|L{40mm}|Y|L{50mm}|}\hline
\hd{항목} & \hd{판정} & \hd{근거(어느 셀)} \\ \hline
\ifwbblank
\rd{9mm}평탄선 높이(\%) &  &  \\ \hline
\rd{9mm}평탄 도달 시점(M) &  &  \\ \hline
\rd{9mm}동일 연령 M4 추세 &  &  \\ \hline
\rd{9mm}판정: 평탄 / 하락 중 / 판독 불가 &  &  \\ \hline
\else
평탄선 높이(\%) & 약 71\% — 27Q1 기준. 26Q4 67\%는 표본 3이라 유보 & 27Q1 M5 = M6 = 10/14 \\ \hline
평탄 도달 시점(M) & M5(27Q1). 26Q4는 M4 & 27Q1 M4 79 → M5 71 → M6 71 \\ \hline
동일 연령 M4 추세 & 개선 중 — 67 → 79 → 83 & 26Q4 2/3 · 27Q1 11/14 · 27Q2 10/12 \\ \hline
판정: 평탄 / 하락 중 / 판독 불가 & 27Q1 기준 M5 이후 71\% 평탄(표본 14). 27Q2는 M3→M4 92→83 하락 중 — 판독 대기. 코호트는 개선 중 & 항목 5가 셋 중 하나의 이름 \\ \hline
\fi
\end{tabularx}}

%% ------------------------------------------------------------------ 4
\landscapeon
\secbar{4. 익명 A·B 비교 — 표본 정의를 물었는가}
\guide{인쇄값을 먼저 적고, 분모 정의를 확인한 뒤 자사 정의로 재계산한 값을 그 아래에 적는다. 두 줄이 같은 회사다.}
{\footnotesize
\begin{tabularx}{\linewidth}{|L{30mm}|Y|L{44mm}|L{44mm}|L{22mm}|}\hline
\hd{회사} & \hd{분모 정의(확인 방법)} & \hd{인쇄값 26Q4/27Q1/27Q2} & \hd{재계산값} & \hd{판독 가능?} \\ \hline
\ifwbblank
\rd{10mm}익명 A — 인쇄 &  &  &  &  \\ \hline
\rd{10mm}익명 A — 자사 정의 재계산 &  &  &  &  \\ \hline
\rd{10mm}익명 B &  &  &  &  \\ \hline
\rd{10mm}자사 &  &  &  &  \\ \hline
\else
익명 A — 인쇄 & 활성 계정(첫 달 말 30일 내 로그인) — 노들에 서면 질의로 확인 & 75 / 93 / 94(분모 4 / 15 / 18) & — & 확인 전: 불가 \\ \hline
익명 A — 자사 정의 재계산 & 총 신규 계약 로고 6 / 20 / 24 & — & 50 / 70 / 71 · M4 67·70·71 & 가능 — 자사보다 낮고 개선 느림 \\ \hline
익명 B & 자사와 동일(확인) & 60 / 59 / 75 & 동일 & 가능 — M4 60·68·75 개선 \\ \hline
자사 & 신규 유료 계약 로고 & 67 / 71 / 83 & — & 가능 \\ \hline
\fi
\end{tabularx}}
\landscapeoff

%% ------------------------------------------------------------------ 5
\secbar{5. 이탈 원인 분류 (5유형)}
\guide{이탈 7건 — 미온보딩 / 챔피언 이탈 / 가치 미실현 / 가격 / 경쟁·대체. 인터뷰 근거.}
{\footnotesize
\begin{tabularx}{\linewidth}{|L{34mm}|L{14mm}|Y|Y|}\hline
\hd{유형} & \hd{건수} & \hd{대표 사례·근거} & \hd{제품 대응} \\ \hline
\ifwbblank
\rd{9mm}미온보딩 &  &  &  \\ \hline
\rd{9mm}챔피언 이탈 &  &  &  \\ \hline
\rd{9mm}가치 미실현 &  &  &  \\ \hline
\rd{9mm}가격 &  &  &  \\ \hline
\rd{9mm}경쟁·대체 &  &  &  \\ \hline
\else
미온보딩 & 2 & 경일정공(27Q1 M2 — 검사원 1명, 첫 달 20건 미만) · 익명 C13(27Q2 M2 — 마스터 미제출 6주) & 첫 달 100건 미만 자동 알림 · 현장 2회 · 마스터 3주 시 계약 보류 \\ \hline
챔피언 이탈 & 1 & 미래오토(27Q1 M3 — 품질팀장 이직, 후임 엑셀 복귀) & 팀장 + 검사원 2인 이상 온보딩 · 단일 챔피언 경고 \\ \hline
가치 미실현 & 2 & 태산정밀(26Q4 M5 — 실사용 2라인, 62항목 부담, 축소 거절 후 해지) · 대양정밀(27Q1 M4 — 배석 공장, 검사원 회피) & 라인 축소 수용(GRR보다 로고) · 검사원 주도 도입 신호를 온보딩 기준에 \\ \hline
가격 & 1 & 신광기공(27Q1 M5 — MES 병행, 프로모션 종료) & 프로모션 종료 60일 전 예고 · MES 병행 고객은 Flow만 \\ \hline
경쟁·대체 & 1 & 익명 C14(27Q2 M3 — 1차사 포털 입력 모듈 개편) & 포털이 못 하는 것(사진·근거·로그)으로 포지셔닝 \\ \hline
\fi
\end{tabularx}}

%% ------------------------------------------------------------------ 6
\secbar{6. 배분 판정서 — 시드 5억을 어디에}
\guide{회사별 배분액·근거·**무배분 사유**. '판독 불가'에 배분하지 않는 것도 판정이다. 자사가 3위라면 그 칸을 채운다.}
{\footnotesize
\begin{tabularx}{\linewidth}{|L{30mm}|L{24mm}|Y|Y|}\hline
\hd{회사} & \hd{배분액} & \hd{근거(어느 표의 어느 셀)} & \hd{무배분 사유 / 조건} \\ \hline
\ifwbblank
\rd{11mm}자사 &  &  &  \\ \hline
\rd{11mm}익명 A &  &  &  \\ \hline
\rd{11mm}익명 B &  &  &  \\ \hline
\rd{11mm}우리 회사가 3위라면 — &  &  &  \\ \hline
\else
자사 & 3.0억 & 항목 3 — 27Q1 M5·M6 71\% 평탄 · M4 67→79→83 · 분모 확인 & 조건: 27Q2 M6 ≥ 70\% 확인 시 집행 \\ \hline
익명 A & 0 & 항목 4 — 인쇄 75/93/94는 활성 계정 분모. 재계산 50/70/71 & 판독 불가 → 무배분(조건부). 총 계약 분모로 재제출 시 재판정. 무배분은 무판정이 아니다 \\ \hline
익명 B & 2.0억 & 항목 4 — 60/59/75, M4 60→68→75 · 분모 동일 & 조건: 27Q1 M6 59\%의 원인 확인 \\ \hline
우리 회사가 3위라면 — & — & 27Q2 M4 83이 M5에서 71로 떨어지면 B와 같은 자리 & 고칠 것: 미온보딩 2 · 가치 미실현 2 — 온보딩 기준·라인 축소 수용 \\ \hline
\fi
\end{tabularx}}

%% ------------------------------------------------------------------ 7
\secbar{7. GRR / NDR 분리}
\guide{이탈만 반영한 GRR(84\%)과 확장 포함 NDR(97\%)을 분리해 적고 분모(기간·기준 MRR)를 명시.}
{\footnotesize
\begin{tabularx}{\linewidth}{|L{24mm}|L{20mm}|Y|Y|}\hline
\hd{지표} & \hd{값} & \hd{분자} & \hd{분모·기간} \\ \hline
\ifwbblank
\rd{9mm}GRR &  &  &  \\ \hline
\rd{9mm}NDR &  &  &  \\ \hline
\rd{9mm}차이의 원천(확장/축소) &  &  &  \\ \hline
\else
GRR & 84\% & 기준 MRR 1,530 − 이탈 245 = 1,285 & D+420 유료 17개사 순 SaaS MRR 1,530 → D+600(6개월) \\ \hline
NDR & 97\% & 1,285 + 확장 199 = 1,484 & 동일. 확장 = 라인 증설 139 + Flow 애드온 60 \\ \hline
차이의 원천(확장/축소) & +13\%p / 0 & NDR < 100 = 순축소 — 확장이 이탈을 못 덮는다 & 12개월 NDR은 없다(회사가 열 달) \\ \hline
\fi
\end{tabularx}}

\end{document}